%%%%%%%%%%%%%%%%%%%%%%%%%%%%%%%%%%%%%%%%%%%%%%%%%%%%%%%%%%%%%%%%%%%%%%
% Lớp 10 | Véctơ (chưa xét tọa độ) | Tích của một số với véctơ
% Dạng 0H5?3-7: Tập hợp điểm
% Số câu: 3  (sinh tự động bởi tach_cau.py; từ nay thầy sửa trực tiếp file này)
%%%%%%%%%%%%%%%%%%%%%%%%%%%%%%%%%%%%%%%%%%%%%%%%%%%%%%%%%%%%%%%%%%%%%%

%%%%%%%%%%%%%%%% TRẮC NGHIỆM 4 PHƯƠNG ÁN %%%%%%%%%%%%%%%%

% nguồn: MaHoaNH_TracNghiem_Lop10
\begin{ex}%[0H5V3-7]%[ma:72fd6de2]
	Cho tam giác $ABC$ cân  tại $C$. Tập hợp các điểm $M$ thỏa mãn đẳng thức $\left|\overrightarrow{MA}+\overrightarrow{MB}\right|=2\left|\overrightarrow{MC}\right|$ là
	\choice 
		{\True Đường thẳng song  song với $AB$}
		{Đường thẳng  vuông góc với $ AB$}
		{Một  điểm}
		{Một  đường tròn}
	 \loigiai{
		Gọi $I$ là trung điểm  $AB$, ta có
		$$\left|\overrightarrow{MA}+\overrightarrow{MB}\right|=2\left|\overrightarrow{MC}\right|\Leftrightarrow 2\left|\overrightarrow{MI}\right|=2\left|\overrightarrow{MC}\right|\Leftrightarrow MI=MC.$$ 
		Điều đó  chứng tỏ điểm $M$ cách đều hai điểm $I$, $C$, nên tập hợp các điểm $M$ là đường trung trực của đoạn $IC$.\\
	 Mặt khác, tam giác $ABC$ cân tại $C$ suy ra $IC\perp AB$ nên đường trung trực của đoạn $IC$ song song với đường thẳng $AB$.}
\end{ex}


%%%%%%%%%%%%%%%% TRẢ LỜI NGẮN %%%%%%%%%%%%%%%%

% nguồn: MaHoaNH_TraLoiNgan_Lop10
\begin{bt}%[0H5H3-7]%[ma:e3ecd340]
	Cho tam giác $ABC$. Gọi $G$ là trọng tâm của $\triangle ABC$. Cho điểm $M$ sao cho $\left|\overrightarrow{MA}+\overrightarrow{MB}+\overrightarrow{MC}\right|=6$,  khi đó điểm $M$ thuộc đường tròn có bán kính bằng bao nhiêu? \shortans{2}
	\loigiai{ 
		Ta có $\overrightarrow{MA}+\overrightarrow{MB}+\overrightarrow{MC}=3 \overrightarrow{M G}\Rightarrow \left|\overrightarrow{MA}+\overrightarrow{MB}+\overrightarrow{MC}\right|=3\left|\overrightarrow{M  G}\right|=3 M G \Leftrightarrow M G=2$.\\
		Vậy tập hợp điểm $M$ là  đường tròn tâm $G$ bán kính bằng $2$.
	} 
\end{bt}

% nguồn: MaHoaNH_TraLoiNgan_Lop10
\begin{bt}%[0H5H3-7]%[ma:97f61da5]
	 Cho tam giác $ABC$ biết $AB=27$ (cm). Điểm $M$ thỏa mãn $\left|2\overrightarrow{MA}+3\overrightarrow{MB}+4\overrightarrow{MC}\right|=\left|\overrightarrow{MB}-\overrightarrow{MA}\right|$. Khi đó điểm $M$ nằm trên đường tròn có bán kính bằng bao nhiêu?
	 \shortans{3}
	 \loigiai{
		 Vì $A$, $B$, $C$ cố định nên ta chọn điểm $I$ thỏa mãn
		\begin{align*}
			 &\,2\overrightarrow{IA}+3\overrightarrow{IB}+4\overrightarrow{IC}=\overrightarrow{0}
			\\ 
			\Leftrightarrow&\, 2\overrightarrow{IA}+3\left(\overrightarrow{IA}+\overrightarrow{IB}\right)+4\left(\overrightarrow{IA}+\overrightarrow{IC}\right)=\overrightarrow{0}\\ \Leftrightarrow&\, 9\overrightarrow{IA}=-3\overrightarrow{AB}-4\overrightarrow{AC} 
			 \\
			\Leftrightarrow&\, \overrightarrow{IA}=-\dfrac{3\overrightarrow{AB}+4\overrightarrow{AC}}{9}. 
		\end{align*}
		Suy ra $I$ duy nhất từ đó $2\overrightarrow{MA}+3\overrightarrow{MB}+4\overrightarrow{MC}=9\overrightarrow{MI}+\left(2\overrightarrow{IA}+3\overrightarrow{IB}+4\overrightarrow{IC}\right)=9\overrightarrow{MI}$  và $\overrightarrow{MA}-\overrightarrow{MB}=\overrightarrow{BA}$.\\
		Từ giả thiết ta suy ra $\left|9\overrightarrow{MI}\right|=\left|\overrightarrow{BA}\right|\Leftrightarrow MI=\dfrac{AB}{9}=3$ (cm). 
	} 
\end{bt}
